\input{configTD}
\usepackage{tikz}
\usepackage{needspace}

% ── Poly à trous : commandes spécifiques ──────────────────────────
\newcommand{\atrou}[2][3cm]{%
  \ifthenelse{\boolean{showSolutions}}{%
    \par\vspace{0.5em}%
    \begin{mdframed}[backgroundcolor=ThemeLight, linewidth=0pt,
      skipabove=0pt, skipbelow=0pt,
      innertopmargin=8pt, innerbottommargin=8pt]%
    #2%
    \end{mdframed}%
    \par\vspace{0.5em}%
  }{%
    \par\vspace{0.5em}%
    \noindent\fbox{\begin{minipage}[c][#1]%
      {\dimexpr\textwidth-2\fboxsep-2\fboxrule\relax}%
      \color{white}\rule{0pt}{0pt}%
    \end{minipage}}%
    \par\vspace{0.5em}%
  }%
}

\newcommand{\val}[1]{%
  \ifthenelse{\boolean{showSolutions}}{#1}{\phantom{#1}}%
}

\newcommand{\R}{\mathbb{R}}
\newcommand{\E}{\mathbb{E}}
\newcommand{\Var}{\operatorname{Var}}
\newcommand{\Xbar}{\overline{X}_n}

\title{\sffamily\bfseries Probabilités --- CM6\\[0.3em]
  \Large Synthèse : choisir le bon outil}
\author{}
\date{}

\begin{document}
\sffamily
\maketitle
\thispagestyle{empty}

\bigskip

\begin{mdframed}[backgroundcolor=ThemeLight, linecolor=Theme,
  linewidth=1.5pt, innertopmargin=10pt, innerbottommargin=10pt]
\textbf{Compétences attendues à l'examen.}

\bigskip

\textit{Événements et probabilités conditionnelles.}
\begin{itemize}\setlength{\itemsep}{0.1em}\setlength{\topsep}{0.1em}
  \item Traduire un énoncé en événements (union, intersection, complémentaire).
  \item Construire et exploiter un arbre pondéré.
  \item Utiliser les probabilités totales et la formule de Bayes.
  \item Caractériser l'indépendance de deux événements ou de deux variables aléatoires.
\end{itemize}

\bigskip

\textit{Lois classiques et calculs associés.}
\begin{itemize}\setlength{\itemsep}{0.1em}\setlength{\topsep}{0.1em}
  \item Reconnaître une loi discrète (Bernoulli, binomiale, géométrique, Poisson) et \emph{justifier ce choix}.
  \item Reconnaître une loi continue (uniforme, exponentielle, normale) à partir de la situation.
  \item Calculer une probabilité par formule de loi, par intégration d'une densité, ou via la fonction de répartition.
  \item Vérifier qu'une fonction est une densité de probabilité.
  \item Calculer l'espérance et la variance d'une variable aléatoire.
  \item Centrer-réduire une loi normale et lire les valeurs de $\Phi$ dans une table donnée.
\end{itemize}

\bigskip

\textit{Du modèle probabiliste au raisonnement statistique.}
\begin{itemize}\setlength{\itemsep}{0.1em}\setlength{\topsep}{0.1em}
  \item Appliquer le théorème central limite pour approcher la loi d'une moyenne ou d'une proportion empirique.
  \item Construire un intervalle de confiance sur une moyenne inconnue.
  \item Choisir une taille d'échantillon pour atteindre une marge d'erreur donnée.
  \item Formuler $H_0$ et $H_1$ en fonction du contexte.
  \item Calculer une statistique de test, une p-value, et conclure au seuil $\alpha$.
  \item Rédiger une conclusion technique cohérente avec le calcul effectué.
\end{itemize}
\end{mdframed}
\newpage


\section*{Sujet d'examen blanc}
% ====================================================================

% --------------------------------------------------------------------
\subsection*{Exercice 1 --- Capteur de contrôle qualité}
% --------------------------------------------------------------------

Sur une chaîne de production, $3\,\%$ des pièces sont défectueuses. Un capteur automatique classe les pièces. Sa fiabilité est la suivante :
\begin{itemize}\setlength{\itemsep}{0.1em}
  \item s'il analyse une pièce défectueuse, il déclenche une alerte dans $92\,\%$ des cas ;
  \item s'il analyse une pièce conforme, il déclenche \emph{quand même} une alerte dans $4\,\%$ des cas (\og fausse alerte \fg).
\end{itemize}

On note $D$ : \og la pièce est défectueuse \fg{} et $S$ : \og le capteur déclenche une alerte \fg.

\textbf{1.} Construire l'arbre pondéré associé.
\atrou[3cm]{%
Premier niveau : $P(D) = 0{,}03$ et $P(\overline{D}) = 0{,}97$.\\
Deuxième niveau : $P(S \mid D) = 0{,}92$, $P(\overline{S} \mid D) = 0{,}08$, $P(S \mid \overline{D}) = 0{,}04$, $P(\overline{S} \mid \overline{D}) = 0{,}96$.
}

\textbf{2.} Calculer $P(S)$.
\atrou[2.5cm]{%
Probabilités totales :
\[
  P(S) = P(S \mid D)P(D) + P(S \mid \overline{D})P(\overline{D})
  = 0{,}92 \times 0{,}03 + 0{,}04 \times 0{,}97 = 0{,}0276 + 0{,}0388 = 0{,}0664.
\]
}

\textbf{3.} Une pièce vient de déclencher une alerte. Quelle est la probabilité qu'elle soit réellement défectueuse\,?
\atrou[3cm]{%
Bayes :
\[
  P(D \mid S) = \dfrac{P(S \mid D)\,P(D)}{P(S)} = \dfrac{0{,}0276}{0{,}0664} \approx 0{,}416.
\]
}

\textbf{4.} Interpréter en une phrase pour un responsable qualité.
\atrou[2cm]{%
Moins d'une alerte sur deux correspond à une vraie pièce défectueuse : il faut prévoir un second contrôle après alerte avant de mettre une pièce au rebut, sinon on jette beaucoup de pièces conformes.
}

% --------------------------------------------------------------------
\needspace{12\baselineskip}
\subsection*{Exercice 2 --- Boulons HR : faut-il financer le procédé préventif\,?}
% --------------------------------------------------------------------

Une entreprise de charpente métallique utilise des boulons haute résistance (HR) livrés par lots de $n = 20$. Avec le procédé de fabrication actuel (\textbf{stratégie A}), chaque boulon est défectueux avec probabilité $p_A = 0{,}05$, indépendamment des autres. Chaque boulon défectueux détecté en atelier coûte $20$~euros de reprise.

Un fournisseur propose un \emph{procédé préventif} (\textbf{stratégie B}) qui abaisse la probabilité de défaut à $p_B = 0{,}01$, mais ajoute $30$~euros de frais fixes par lot.

\medskip

\textit{Données utiles (sans calculatrice)~:} $0{,}95^{20} \approx 0{,}36$ \quad et \quad $0{,}99^{20} \approx 0{,}82$.

\textbf{1.} Soit $X$ le nombre de boulons défectueux dans un lot. Déterminer la loi de $X$ et préciser ses paramètres dans chacune des deux stratégies.
\atrou[2.5cm]{%

$n=20$ essais indépendants, chacun \og défectueux ou non \fg{} avec une probabilité constante : c'est exactement le cadre de la loi binomiale.
\[
  \text{Stratégie A :}\quad X_A \sim \mathcal{B}(20,\,0{,}05),
  \qquad
  \text{Stratégie B :}\quad X_B \sim \mathcal{B}(20,\,0{,}01).
\]
}

\textbf{2.} Calculer le coût moyen total $\E(C_A)$ d'un lot sous la stratégie A.
\atrou[2.5cm]{%
$\E(X_A) = n p_A = 20 \times 0{,}05 = 1$. \\
Coût moyen de reprise : $\E(C_A) = 20 \times \E(X_A) = 20 \times 1 = 20$~euros par lot.
}

\textbf{3.} Calculer le coût moyen total $\E(C_B)$ d'un lot sous la stratégie B.
\atrou[2.5cm]{%
$\E(X_B) = 20 \times 0{,}01 = 0{,}2$. \\
Coût moyen : $\E(C_B) = 30 + 20 \times \E(X_B) = 30 + 20 \times 0{,}2 = 30 + 4 = 34$~euros par lot.
}

\textbf{4.} Comparer la probabilité qu'un lot contienne \emph{au moins un} boulon défectueux sous chacune des deux stratégies.
\atrou[3cm]{%
\[
  P(X_A \ge 1) = 1 - P(X_A = 0) = 1 - 0{,}95^{20} \approx 1 - 0{,}36 = 0{,}64.
\]
\[
  P(X_B \ge 1) = 1 - P(X_B = 0) = 1 - 0{,}99^{20} \approx 1 - 0{,}82 = 0{,}18.
\]
La prévention divise le risque \og au moins un défaut \fg{} par plus de $3$ ($64\,\%$ contre $18\,\%$).
}

\textbf{5.} Quelle stratégie recommanderiez-vous à un client \emph{très sensible au risque} de non-conformité (par exemple un ouvrage de sécurité)\,? Justifier en mobilisant les deux comparaisons précédentes.
\atrou[3cm]{%
La stratégie B coûte $14$~euros de plus par lot en moyenne ($34$ vs $20$), mais réduit fortement le risque qu'un lot contienne au moins un défaut ($18\,\%$ vs $64\,\%$). Pour un client sensible au risque (ouvrage critique, conséquences lourdes en cas de défaillance), le surcoût est justifié et on recommande la stratégie B. Pour un client peu exposé, la stratégie A est plus économique en moyenne. \emph{Le critère du coût moyen ne suffit pas seul à trancher~: il faut aussi tenir compte de la queue de distribution.}
}

% --------------------------------------------------------------------
\needspace{12\baselineskip}
\subsection*{Exercice 3 --- Épaisseur d'enrobé routier}
% --------------------------------------------------------------------

L'épaisseur $E$ d'une couche d'enrobé bitumineux mesurée en un point d'une chaussée suit une loi normale de moyenne $8$~cm et d'écart-type $0{,}5$~cm. Le cahier des charges impose une épaisseur d'au moins $7$~cm en chaque point contrôlé.

\textbf{1.} Quelle est la probabilité qu'un point contrôlé soit \emph{non conforme}\,?
\atrou[3cm]{%
On standardise : $Z = (E - 8)/0{,}5$.
\[
  P(E < 7) = P\!\left(Z < \dfrac{7-8}{0{,}5}\right) = P(Z < -2) = 1 - \Phi(2) \approx 1 - 0{,}9772 = 0{,}0228.
\]
Environ $2{,}3\,\%$ des points contrôlés ne respectent pas le cahier des charges.
}

\textbf{2.} Quelle est la probabilité que l'épaisseur soit comprise entre $7{,}5$ et $8{,}5$~cm\,?
\atrou[2.5cm]{%
\[
  P(7{,}5 \le E \le 8{,}5) = P(-1 \le Z \le 1) = 2\Phi(1) - 1 \approx 2 \times 0{,}8413 - 1 = 0{,}6826.
\]
}

\textbf{3.} Déterminer le seuil $e_0$ tel que $90\,\%$ des points contrôlés aient une épaisseur supérieure à $e_0$.
\atrou[3cm]{%
On cherche $e_0$ tel que $P(E > e_0) = 0{,}9$, soit $P(E \le e_0) = 0{,}1$. On lit $\Phi^{-1}(0{,}1) = -1{,}28$ (par symétrie). Donc
\[
  e_0 = 8 + 0{,}5 \times (-1{,}28) = 7{,}36 \text{ cm}.
\]
}

\textbf{4.} En une phrase, ce que représente $e_0$ pour le contrôle qualité.
\atrou[2cm]{%
$e_0 = 7{,}36$~cm est l'épaisseur \og garantie à $90\,\%$ \fg{} : neuf points contrôlés sur dix dépassent cette valeur, ce qui permet de dimensionner le risque d'un sondage de réception.
}

% --------------------------------------------------------------------
\needspace{12\baselineskip}
\subsection*{Exercice 4 --- Contrôle d'un lot de ferraillages}
% --------------------------------------------------------------------

Une entreprise de préfabrication contrôle ses ferraillages à la sortie de la chaîne de soudure. D'après ses statistiques internes, le taux de défauts est annoncé à $5\,\%$. Sur un échantillon de $n = 400$ ferraillages prélevés au hasard dans la production d'une journée, on observe $36$ défauts, soit une proportion empirique $\widehat{p} = 0{,}09$.

\textbf{1.} Justifier qu'on peut utiliser une approximation normale pour la proportion empirique $\widehat{p}$.
\atrou[2.5cm]{%
On vérifie les conditions usuelles : $n\widehat{p} = 36 \ge 10$ et $n(1-\widehat{p}) = 364 \ge 10$. Par le TCL,
\[
  \widehat{p} \approx \mathcal{N}\!\left(p,\;\dfrac{p(1-p)}{n}\right).
\]
}

\textbf{2.} Construire un intervalle de confiance à $95\,\%$ pour la vraie proportion $p$ de ferraillages défectueux.
\atrou[4cm]{%
On remplace $p$ inconnu par $\widehat{p}$ dans l'écart-type :
\[
  IC_{95\%}(p) = \widehat{p} \pm 1{,}96\,\sqrt{\dfrac{\widehat{p}(1-\widehat{p})}{n}}
  = 0{,}09 \pm 1{,}96\,\sqrt{\dfrac{0{,}09 \times 0{,}91}{400}}.
\]
\[
  \sqrt{\dfrac{0{,}09 \times 0{,}91}{400}} = \sqrt{2{,}048 \cdot 10^{-4}} \approx 0{,}0143.
\]
\[
  IC_{95\%}(p) \approx 0{,}09 \pm 0{,}028 = [0{,}062\,;\,0{,}118].
\]
}

\textbf{3.} L'annonce \og taux de défauts $= 5\,\%$ \fg{} est-elle compatible avec ces observations\,? Conclure en utilisant la dualité IC $\Leftrightarrow$ test.
\atrou[2.5cm]{%
$p_0 = 0{,}05$ est \emph{en dehors} de l'IC à $95\,\%$ (à gauche). Par dualité, le test bilatéral $H_0 : p = 0{,}05$ contre $H_1 : p \neq 0{,}05$ est \textbf{rejeté} au seuil $5\,\%$ : l'annonce est incompatible avec les données.
}

\textbf{4.} Mettre en place un test d'hypothèse formel. La question pratique est~: \og le taux dépasse-t-il $5\,\%$\,? \fg{}. Préciser $H_0$ et $H_1$, calculer $Z_{\text{obs}}$ sous $H_0$ et la p-value, puis conclure au seuil $\alpha = 5\,\%$.
\atrou[4cm]{%
Le client cherche à attraper un dépassement, donc test \textbf{unilatéral à droite} :
\[
  H_0 : p \le 0{,}05 \quad \text{contre} \quad H_1 : p > 0{,}05.
\]
\emph{Attention~:} sous $H_0$, on utilise $p_0 = 0{,}05$ (cas limite) pour calculer l'écart-type, pas $\widehat{p}$.
\[
  \sigma_0 = \sqrt{\dfrac{0{,}05 \times 0{,}95}{400}} = \sqrt{1{,}19 \cdot 10^{-4}} \approx 0{,}0109.
\]
\[
  Z_{\text{obs}} = \dfrac{\widehat{p} - 0{,}05}{\sigma_0} = \dfrac{0{,}09 - 0{,}05}{0{,}0109} \approx 3{,}67.
\]
\[
  \text{p-value} = P(Z \ge 3{,}67 \mid H_0) = 1 - \Phi(3{,}67) \approx 1{,}2 \cdot 10^{-4}.
\]
$1{,}2 \cdot 10^{-4} \ll 0{,}05$ : on \textbf{rejette} $H_0$.
}

\textbf{5.} Rédiger une conclusion technique courte.
\atrou[2cm]{%
Les observations sont statistiquement incompatibles avec un taux de défauts de $5\,\%$ : le vrai taux est très probablement supérieur. Il est justifié d'alerter le service qualité et de relancer un contrôle complet de la chaîne de soudure avant livraison.
}

\vspace{1em}

\newpage

\section*{Table de la loi normale centrée réduite}

\vspace{1cm}

\begin{center}
    \begin{tabular}{|r|r@{\ }r@{\ }r@{\ }r@{\ }r@{\ }r@{\ }r@{\ }r@{\ }r@{\ }r@{\ }|}
    \hline
    \multicolumn{11}{|c|}{ Fonction de répartition de la loi $\mathcal{N}(0,1)$}\\
    \hline
    \hline
    $x$&0.00&0.01&0.02&0.03&0.04&0.05&0.06&0.07&0.08&0.09\\
    \hline
    0.0&0.5000&0.5040&0.5080&0.5120&0.5160&0.5199&0.5239&0.5279&0.5319&0.5359\\
    \hline
    0.1&0.5398&0.5438&0.5478&0.5517&0.5557&0.5596&0.5636&0.5675&0.5714&0.5753\\
    \hline
    0.2&0.5793&0.5832&0.5871&0.5910&0.5948&0.5987&0.6026&0.6064&0.6103&0.6141\\
    \hline
    0.3&0.6179&0.6217&0.6255&0.6293&0.6331&0.6368&0.6406&0.6443&0.6480&0.6517\\
    \hline
    0.4&0.6554&0.6591&0.6628&0.6664&0.6700&0.6736&0.6772&0.6808&0.6844&0.6879\\
    \hline
    0.5&0.6915&0.6950&0.6985&0.7019&0.7054&0.7088&0.7123&0.7157&0.7190&0.7224\\
    \hline
    0.6&0.7257&0.7291&0.7324&0.7357&0.7389&0.7422&0.7454&0.7486&0.7517&0.7549\\
    \hline
    0.7&0.7580&0.7611&0.7642&0.7673&0.7703&0.7734&0.7764&0.7794&0.7823&0.7852\\
    \hline
    0.8&0.7881&0.7910&0.7939&0.7967&0.7995&0.8023&0.8051&0.8078&0.8106&0.8133\\
    \hline
    0.9&0.8159&0.8186&0.8212&0.8238&0.8264&0.8289&0.8315&0.8340&0.8365&0.8389\\
    \hline
    1.0&0.8413&0.8438&0.8461&0.8485&0.8508&0.8531&0.8554&0.8577&0.8599&0.8621\\
    \hline
    1.1&0.8643&0.8665&0.8686&0.8708&0.8729&0.8749&0.8770&0.8790&0.8810&0.8830\\
    \hline
    1.2&0.8849&0.8869&0.8888&0.8907&0.8925&0.8944&0.8962&0.8980&0.8997&0.9015\\
    \hline
    1.3&0.9032&0.9049&0.9066&0.9082&0.9099&0.9115&0.9131&0.9147&0.9162&0.9177\\
    \hline
    1.4&0.9192&0.9207&0.9222&0.9236&0.9251&0.9265&0.9279&0.9292&0.9306&0.9319\\
    \hline
    1.5&0.9332&0.9345&0.9357&0.9370&0.9382&0.9394&0.9406&0.9418&0.9429&0.9441\\
    \hline
    1.6&0.9452&0.9463&0.9474&0.9484&0.9495&0.9505&0.9515&0.9525&0.9535&0.9545\\
    \hline
    1.7&0.9554&0.9564&0.9573&0.9582&0.9591&0.9599&0.9608&0.9616&0.9625&0.9633\\
    \hline
    1.8&0.9641&0.9649&0.9656&0.9664&0.9671&0.9678&0.9686&0.9693&0.9699&0.9706\\
    \hline
    1.9&0.9713&0.9719&0.9726&0.9732&0.9738&0.9744&0.9750&0.9756&0.9761&0.9767\\
    \hline
    2.0&0.9772&0.9778&0.9783&0.9788&0.9793&0.9798&0.9803&0.9808&0.9812&0.9817\\
    \hline
    2.1&0.9821&0.9826&0.9830&0.9834&0.9838&0.9842&0.9846&0.9850&0.9854&0.9857\\
    \hline
    2.2&0.9861&0.9864&0.9868&0.9871&0.9875&0.9878&0.9881&0.9884&0.9887&0.9890\\
    \hline
    2.3&0.9893&0.9896&0.9898&0.9901&0.9904&0.9906&0.9909&0.9911&0.9913&0.9916\\
    \hline
    2.4&0.9918&0.9920&0.9922&0.9925&0.9927&0.9929&0.9931&0.9932&0.9934&0.9936\\
    \hline
    2.5&0.9938&0.9940&0.9941&0.9943&0.9945&0.9946&0.9948&0.9949&0.9951&0.9952\\
    \hline
    2.6&0.9953&0.9955&0.9956&0.9957&0.9959&0.9960&0.9961&0.9962&0.9963&0.9964\\
    \hline
    2.7&0.9965&0.9966&0.9967&0.9968&0.9969&0.9970&0.9971&0.9972&0.9973&0.9974\\
    \hline
    2.8&0.9974&0.9975&0.9976&0.9977&0.9977&0.9978&0.9979&0.9979&0.9980&0.9981\\
    \hline
    2.9&0.9981&0.9982&0.9982&0.9983&0.9984&0.9984&0.9985&0.9985&0.9986&0.9986\\
    \hline
    3.0&0.9987&0.9987&0.9987&0.9988&0.9988&0.9989&0.9989&0.9989&0.9990&0.9990\\
    \hline
    3.1&0.9990&0.9991&0.9991&0.9991&0.9992&0.9992&0.9992&0.9992&0.9993&0.9993\\
    \hline
    3.2&0.9993&0.9993&0.9994&0.9994&0.9994&0.9994&0.9994&0.9995&0.9995&0.9995\\
    \hline
    3.3&0.9995&0.9995&0.9995&0.9996&0.9996&0.9996&0.9996&0.9996&0.9996&0.9997\\
    \hline
    3.4&0.9997&0.9997&0.9997&0.9997&0.9997&0.9997&0.9997&0.9997&0.9997&0.9998\\
    \hline
    3.5&0.9998&0.9998&0.9998&0.9998&0.9998&0.9998&0.9998&0.9998&0.9998&0.9998\\
    \hline
    3.6&0.9998&0.9999&0.9999&0.9999&0.9999&0.9999&0.9999&0.9999&0.9999&0.9999\\
    \hline
    3.7&0.9999&0.9999&0.9999&0.9999&0.9999&0.9999&0.9999&0.9999&0.9999&0.9999\\
    \hline
    3.8&0.9999&0.9999&0.9999&0.9999&0.9999&0.9999&0.9999&0.9999&0.9999&0.9999\\
    \hline
    3.9&1.0000&1.0000&1.0000&1.0000&1.0000&1.0000&1.0000&1.0000&1.0000&1.0000\\
    \hline
    \end{tabular}
\end{center}

\end{document}
